\chapter{Fichiers, commandes et reproduction}
\section{Installer le projet et vérifier son fonctionnement}
Les commandes suivantes sont à lancer sous Linux depuis la racine du
dépôt. Python 3.10 ou plus récent est déclaré compatible. Les contrôles
décrits dans le rapport ont été exécutés sous Python 3.12. La matrice CI
du dépôt prévoit Python 3.10 et 3.12 ; le rapport ne suppose pas que
l'exécution distante de cette CI a déjà été observée.
\begin{commandes}
git clone https://github.com/birane-m/ring-star-problem.git
cd ring-star-problem
make install
make test
\end{commandes}
Les dépendances sont définies dans \code{pyproject.toml}. Le fichier
\code{requirements.txt} installe le projet en mode éditable et ne
constitue pas une seconde liste indépendante. Les versions de l'exécution
archivée se trouvent dans \code{resultats/campagne/environment.json}.

\section{Reproduire une résolution et une campagne}
L'exemple individuel att48 est paramétré par :
\begin{commandes}
make tabu INSTANCE=data/instances/att48.tsp P=10 ALPHA=1 \
  SEED=0 ITERATIONS=100 CANDIDATES=15 TENURE=10
\end{commandes}
L'interface explicite permet aussi de fixer tous les paramètres :
\begin{commandes}
.venv/bin/ring-star run tabu data/instances/att48.tsp 10 \
  --alpha 1 --seed 0 --iterations 100 \
  --candidates-per-iteration 15 --tabu-tenure 10
\end{commandes}
La campagne de démonstration se lance avec :
\begin{commandes}
make benchmark
\end{commandes}
Elle lit \code{configs/benchmark_demo.json}. Une relance crée un nouveau
dossier horodaté. Les durées, et éventuellement le comportement de CBC
au voisinage du temps limite, peuvent varier. Pour retrouver exactement
les chiffres de ce rapport, utiliser les résultats archivés plutôt que
supposer l'identité des durées d'une nouvelle exécution.

\section{Retrouver les sorties}
\begin{table}[htbp]\small
\begin{tabularx}{\textwidth}{lL}\toprule
Fichier ou dossier & Rôle \\\midrule
\code{outputs/index.md} & Point d'entrée vers les exécutions et campagnes. \\
\code{summary.md} & Lecture d'une résolution, de ses coûts et de ses figures. \\
\code{solution.json} & Coordonnées, solution, paramètres, coûts et statut. \\
\code{initial_solution.json} & Référence initiale d'une recherche. \\
\code{history.csv} & Coûts courant et meilleur, évaluations et temps de recherche. \\
\code{runs.csv} & Enregistrements détaillés des méthodes et graines d'une campagne. \\
\code{summary.csv} & Moyennes, dispersions, statuts et liens aux meilleures graines. \\
\code{completion.json} & État terminé ou interrompu et compte des enregistrements. \\
\code{environment.json} & Versions et empreinte des sources. \\
\code{figures/} & Figures de solution, convergence et comparaison. \\
\bottomrule\end{tabularx}
\caption{Repères de navigation dans les sorties du projet.}
\end{table}
Les anciens chemins absolus présents dans certains fichiers de provenance
documentent leur origine. Les solutions contiennent les coordonnées et
les chemins utiles à la lecture du rapport sont relatifs à ses sources.
La régénération des tableaux ne dépend pas des anciens chemins absolus.

\section{Compiler le rapport et régénérer les tableaux}
L'archive des sources est autonome pour la compilation ; elle contient
les figures utilisées, les tableaux et les résultats qui les alimentent.
Les paquets LaTeX usuels de mathématiques, de langue française et de
mise en page sont nécessaires, ainsi que Python 3 pour les contrôles.
Depuis le dossier contenant \code{main.tex} :
\begin{commandes}
python3 scripts/generer_tableaux.py
latexmk -pdf -interaction=nonstopmode -halt-on-error main.tex
\end{commandes}
Sans \code{latexmk}, compiler \code{main.tex} trois fois avec pdfLaTeX.
La première commande reconstruit les tableaux, vérifie les coûts des
172 solutions et documente le rapprochement des 191 lignes originales
avec les 192 cas attendus. Elle ne lance aucun solveur.

La couverture conserve l'image fournie. Son URL source et ses crédits
sont rassemblés dans \code{IMAGE_SOURCE.txt}. Les figures de résultats
reprennent les exports vectoriels de la campagne et de l'exemple individuel.

\section{Intégrer le correctif d'export au dépôt}
Le fichier \code{correctifs/export_cas_ignores.patch} contient uniquement
l'écriture finale du CSV et le renforcement du test existant. Après avoir
placé les sources du rapport dans \code{reports/latex_ring_star/},
on peut l'appliquer depuis la racine du dépôt :
\begin{commandes}
git apply --check \
  reports/latex_ring_star/correctifs/export_cas_ignores.patch
git apply \
  reports/latex_ring_star/correctifs/export_cas_ignores.patch
make test
\end{commandes}
Le correctif reste à intégrer dans un commit distinct ; l'export original
archivé est conservé.
